%*******************************************************
% How to read this book
%*******************************************************
\pdfbookmark[1]{How to Read This Book}{howtoread}
\chapter*{How to Read This Book}
\markboth{\spacedlowsmallcaps{How to Read This Book}}{\spacedlowsmallcaps{How to Read This Book}}

\section*{What you need}
If you remember a little algebra from school (how to multiply out brackets) and the basic idea of sine and cosine, you have everything you need. Each new idea is introduced slowly, first with a picture, then in words, and only then with symbols.

\section*{Two ways through the book}
\begin{description}
  \item[The short road.] Read Chapters~1 and~2, then Chapter~4, then Chapter~6, and finish with the Coda. This takes you straight to the heart of the theory. A short reminder of the General Sinus is given where Chapter~4 needs it, and \Cref{app:general-sinus} lists its formulas. To see where the theory leads, read Chapter~9 before the Coda.
  \item[The full road.] Read everything in order. Chapter~3 shows where the General Sinus comes from, and Chapter~5 explains why physics sometimes needs its language to grow. Chapter~7 draws all the particles, ions and nuclei in one frame, and Part~\ref{pt:nuclei-fields} follows the cosmic mass into the nuclei (Chapter~8) and into bodies in motion and their fields (Chapter~9). Part~\ref{pt:open} ends the road with an open discussion of the Coupling Force (Chapter~10): what holds, what is excluded, and what is still open; and with the whole frame (Chapter~11): antimatter, the half of negative mass, and the black holes.
\end{description}

\section*{What the boxes mean}
\begin{plainwords}
Boxes like this one explain an idea in everyday language. You can read them before or after the formal statement they accompany. They never introduce new results.
\end{plainwords}

\begin{workedexample}
Boxes like this one do a short calculation, step by step, using the formulas exactly as they are stated in the text.
\end{workedexample}

A line in small capitals, such as

\scene{Paris, 1785}

\noindent opens a scene from history: a person, a place, and a moment when an idea changed.

\section*{The formal statements}
The theory itself is stated in numbered \emph{axioms}, \emph{definitions}, \emph{theorems} and \emph{remarks}, and, in Chapters~8 and~9, \emph{laws}: statements that are not proved but tested against measurement. Chapter~10 adds \emph{views}: readings put forward for discussion, neither proved nor yet tested, each with its status.

\section*{Symbols}
A table of all the symbols, units and constants used in the book is given in \Cref{app:notation}.
